\chapter{المكعبات والصناديق}\label{ch:g5:solids}

لنغادر الصفحة المستوية: \omterm{def:g5:solids:def}{المجسّمات} هي أشكال العالم الحقيقي ---
حجر النرد، وعلب الحبوب، والعلب المعدنية، والكرات. يتعلم هذا الفصل
تسميتها، وعدّ أوجهها وأحرفها، ونشرها تصاميمَ مستوية، وقياس
حجمها بعدّ المكعبات الصغيرة. (وصيغ \omterm{def:g5:solids:volume}{الحجم} تأتي في
\cref{ch:g6:measure}.)

\section{عائلة المجسّمات}

\begin{definition}[المجسّمات ومفرداتها]\label{def:g5:solids:def}
\emph{المجسّم}\index{مجسم} يشغل حيزًا. وجوانبه المستوية
\emph{أوجه}\index{وجه}؛ وتلتقي الأوجه على \emph{أحرف}\index{حرف}؛
وتلتقي الأحرف في \emph{رؤوس}\index{رأس}. وصور
العائلة:
\begin{itemize}
  \item \emph{المكعب}\index{مكعب}: $6$ أوجه مربعة، و $12$ حرفًا،
        و $8$ رؤوس؛
  \item و\emph{الصندوق}\index{صندوق} (متوازي المستطيلات): $6$
        أوجه مستطيلة، و $12$ حرفًا، و $8$ رؤوس؛
  \item و\emph{الأسطوانة}\index{أسطوانة}: وجهان قرصيان ووجه ملفوف
        (علبة معدنية)؛
  \item و\emph{الهرم}\index{هرم}: قاعدة مضلعة ومثلثات تلتقي في
        رأس؛
  \item و\emph{الكرة}: سطح مستدير تمامًا، لا وجه له ولا
        حرف.
\end{itemize}
\end{definition}

\begin{omfigure}
\begin{tikzpicture}[scale=0.62]
  % cube
  \draw[thick] (0,0) rectangle (1.8,1.8);
  \draw[thick] (0,1.8) -- (0.7,2.5) -- (2.5,2.5) -- (1.8,1.8);
  \draw[thick] (2.5,2.5) -- (2.5,0.7) -- (1.8,0);
  \node[below, font=\small] at (1.2,-0.3) {مكعب};
  % box
  \begin{scope}[xshift=4.4cm]
  \draw[thick] (0,0) rectangle (2.8,1.4);
  \draw[thick] (0,1.4) -- (0.6,2.0) -- (3.4,2.0) -- (2.8,1.4);
  \draw[thick] (3.4,2.0) -- (3.4,0.6) -- (2.8,0);
  \node[below, font=\small] at (1.7,-0.3) {صندوق};
  \end{scope}
  % cylinder
  \begin{scope}[xshift=9.6cm]
  \draw[thick] (0,0.35) ellipse (0.85 and 0.3);
  \draw[thick] (-0.85,0.35) -- (-0.85,2.2);
  \draw[thick] (0.85,0.35) -- (0.85,2.2);
  \draw[thick] (0,2.2) ellipse (0.85 and 0.3);
  \node[below, font=\small] at (0,-0.3) {أسطوانة};
  \end{scope}
  % pyramid
  \begin{scope}[xshift=12.6cm]
  \draw[thick] (-0.9,0.2) -- (0.9,0.2) -- (1.35,0.75) -- cycle;
  \draw[thick] (-0.9,0.2) -- (0.15,2.4) -- (0.9,0.2);
  \draw[thick] (1.35,0.75) -- (0.15,2.4);
  \draw[dashed] (-0.9,0.2) -- (-0.45,0.75) -- (1.35,0.75);
  \draw[dashed] (-0.45,0.75) -- (0.15,2.4);
  \node[below, font=\small] at (0.2,-0.3) {هرم};
  \end{scope}
  % ball
  \begin{scope}[xshift=16.4cm]
  \draw[thick] (0,1.2) circle (1.1);
  \draw[dashed] (0,1.2) ellipse (1.1 and 0.35);
  \node[below, font=\small] at (0,-0.3) {كرة};
  \end{scope}
\end{tikzpicture}

{\small معرض \omterm{def:g5:solids:def}{المجسّمات}. والخطوط المتقطعة تُظهر \omterm{def:g5:solids:def}{الأحرف} المخفية ---
ظهر \omterm{def:g5:solids:def}{المجسّم}، مرئيًّا من خلاله.}
\end{omfigure}

\begin{example}[العدّ على المكعب]\label{ex:g5:solids:count}
تحقّق على حجر نرد: $6$ \omterm{def:g5:solids:def}{أوجه} (الأعداد $1$ إلى $6$)، و $8$ \omterm{def:g5:solids:def}{رؤوس}
(الأركان)، و $12$ حرفًا (مرِّر إصبعك عليها: $4$ في الأعلى، و $4$
في الأسفل، و $4$ رأسية). و
$6 + 8 = 12 + 2$ --- والنمط الغريب نفسه يصحّ للصندوق
وللهرم (و\cref{ex:g8:solids:count} تعود إليه).
\end{example}

\section{التصاميم}

\begin{definition}[التصميم]\label{def:g5:solids:net}
\emph{تصميم}\index{تصميم} \omterm{def:g5:solids:def}{المجسّم} رسم مستوٍ لكل
أوجهه، متصلة على \omterm{def:g5:solids:def}{الأحرف}، ينطوي فيعطي \omterm{def:g5:solids:def}{المجسّم} --- أي
\omterm{def:g5:solids:def}{المجسّم} مفرودًا كصندوق كرتوني مبسوط لإعادة التدوير.
\end{definition}

\begin{omfigure}
\begin{tikzpicture}[scale=0.5]
  \draw[gray!40, thin] (0,0) grid (4,3);
  \draw[thick, omDef] (0,1) rectangle (1,2);
  \draw[thick, omDef] (1,1) rectangle (2,2);
  \draw[thick, omDef] (2,1) rectangle (3,2);
  \draw[thick, omDef] (3,1) rectangle (4,2);
  \draw[thick, omDef] (1,2) rectangle (2,3);
  \draw[thick, omDef] (1,0) rectangle (2,1);
  \node[below, font=\small] at (2,-0.5) {تصميم للمكعب};
  \begin{scope}[xshift=7.5cm]
  \draw[gray!40, thin] (0,0) grid (4,3);
  \draw[thick, omThm] (0,1) rectangle (1,2);
  \draw[thick, omThm] (1,1) rectangle (2,2);
  \draw[thick, omThm] (2,1) rectangle (3,2);
  \draw[thick, omThm] (3,1) rectangle (4,2);
  \draw[thick, omThm] (0,2) rectangle (1,3);
  \draw[thick, omThm] (0,0) rectangle (1,1);
  \node[below, font=\small] at (2,-0.5) {تصميم صحيح آخر};
  \end{scope}
  \begin{scope}[xshift=15cm]
  \draw[gray!40, thin] (0,0) grid (4,3);
  \draw[thick, omExo] (0,1) rectangle (1,2);
  \draw[thick, omExo] (1,1) rectangle (2,2);
  \draw[thick, omExo] (2,1) rectangle (3,2);
  \draw[thick, omExo] (3,1) rectangle (4,2);
  \draw[thick, omExo] (0,2) rectangle (1,3);
  \draw[thick, omExo] (1,2) rectangle (2,3);
  \node[below, font=\small] at (2,-0.5) {ليس تصميمًا! (لماذا؟)};
  \end{scope}
\end{tikzpicture}

{\small ثلاثة ترتيبات لستة مربعات. الأولان ينطويان فيعطيان
مكعبًا؛ أما الثالث فلا --- فعند الطيّ يقع مربعان على
\omterm{def:g5:solids:def}{الوجه} نفسه ويبقى \omterm{def:g5:solids:def}{المكعب} مفتوحًا. قصّها وجرّب!}
\end{omfigure}

\begin{example}[قراءة تصميم]\label{ex:g5:solids:read}
في حجر النرد، \omterm{def:g1:addition:def}{مجموع} الوجهين المتقابلين $7$. وفي \omterm{def:g5:solids:net}{التصميم}، الوجهان
المتقابلان هما اللذان يفصل بينهما مربع واحد بالضبط في صفّ (أو حول
ركن) --- لا الجاران! وتعليم الأزواج على \omterm{def:g5:solids:net}{التصميم} قبل
الطيّ تمرين ممتاز في رؤية الفضاء مستويًا.
\end{example}

\section{الحجم: عدّ المكعبات}

\begin{definition}[الحجم بالعدّ]\label{def:g5:solids:volume}
\emph{حجم}\index{حجم} \omterm{def:g5:solids:def}{مجسّم} مبنيّ من مكعبات واحدية هو
عدد المكعبات التي يحتويها. والوحدة قد تكون \omterm{def:g5:solids:def}{المكعب} السنتيمتري
(سم$^3$) أو \omterm{def:g5:solids:def}{المكعب} المتري (م$^3$).
\end{definition}

\begin{example}[عدّ صندوق طبقةً طبقة]\label{ex:g5:solids:layers}
صندوق طوله $4$ مكعبات، وعرضه $3$، وارتفاعه $2$: الطبقة السفلى
تسع $4 \times 3 = 12$ مكعبًا، والطبقات $2$:
\[
12 \times 2 = 24 \text{ مكعبًا} .
\]
والعدّ بالطبقات ينجح دائمًا --- وهو يبرهن ضمنًا الصيغة
$V = L \times w \times h$ في \cref{ch:g6:measure}.
\end{example}

\begin{omfigure}
\begin{tikzpicture}[scale=0.52]
  \foreach \y in {0,1} {
    \foreach \x in {0,...,3} {
      \foreach \z in {0,1,2} {
        \begin{scope}[shift={({\x+0.4*\z},{\y+0.28*\z})}]
          \draw[fill=omDef!12] (0,0) rectangle (1,1);
          \draw[fill=omDef!22] (0,1) -- (0.4,1.28) -- (1.4,1.28)
            -- (1,1) -- cycle;
          \draw[fill=omDef!30] (1,0) -- (1.4,0.28) -- (1.4,1.28)
            -- (1,1) -- cycle;
        \end{scope}
      }
    }
  }
  \node[below, font=\small] at (2.6,-0.4)
    {$4 \times 3$ مكعبًا في الطبقة، و $2$ من الطبقات: $24$ مكعبًا};
\end{tikzpicture}

{\small صندوق $4 \times 3 \times 2$ من المكعبات الواحدية، معدودًا
طبقةً طبقة.}
\end{omfigure}

\begin{example}[المكعبات نفسها، وأشكال مختلفة]\label{ex:g5:solids:rearrange}
ثمانية مكعبات واحدية يمكن أن تبني مكعبًا $2 \times 2 \times 2$، أو
عصا $8 \times 1 \times 1$، أو لوحًا $4 \times 2 \times 1$: ثلاثة
\omterm{def:g5:solids:def}{مجسّمات} مختلفة، \omterm{def:g5:solids:volume}{وحجم} واحد ($8$ مكعبات). \omterm{def:g5:solids:volume}{فالحجم} يعدّ
المادة، لا الشكل.
\end{example}

\section{تمارين}

\begin{exercise}[$\star$]\label{exo:g5:solids:1}
سمِّ شيئًا حقيقيًّا على شكل: \omterm{def:g5:solids:def}{مكعب}؛ وصندوق؛ \omterm{def:g5:solids:def}{وأسطوانة}؛ وكرة؛
\omterm{def:g5:solids:def}{وهرم}.
\end{exercise}

\begin{exercise}[$\star$]\label{exo:g5:solids:2}
كم وجهًا وحرفًا ورأسًا للصندوق؟ وتحقّق من النمط
\omterm{def:g5:solids:def}{الأوجه} $+$ \omterm{def:g5:solids:def}{الرؤوس} $=$ \omterm{def:g5:solids:def}{الأحرف} $+ 2$.
\end{exercise}

\begin{exercise}[$\star$]\label{exo:g5:solids:3}
أي \omterm{def:g5:solids:def}{أوجه} الصندوق متطابقة؟ (اجمع \omterm{def:g5:solids:def}{الأوجه} الستة أزواجًا.)
وما الخاصّ في \omterm{def:g5:solids:def}{أوجه} \omterm{def:g5:solids:def}{المكعب}؟
\end{exercise}

\begin{exercise}[$\star$]\label{exo:g5:solids:4}
اُرسم تصميمًا \omterm{def:g5:solids:def}{لمكعب} طول ضلعه $2$ سم (اِستعمل النموذج المتصالب
في الفصل)، وقصّه واطوِه.
\end{exercise}

\begin{exercise}[$\star$]\label{exo:g5:solids:5}
على \omterm{def:g5:solids:net}{التصميم} المتصالب لحجر نرد، ضع الأعداد $1$ إلى $6$ بحيث
يكون \omterm{def:g1:addition:def}{مجموع} الوجهين المتقابلين $7$. (واِستعمل \cref{ex:g5:solids:read}
لإيجاد الأزواج المتقابلة أولًا.)
\end{exercise}

\begin{exercise}[$\star$]\label{exo:g5:solids:6}
عُدّ المكعبات الواحدية: صندوق طوله $5$، وعرضه $2$، وارتفاعه $3$؛
\omterm{def:g5:solids:def}{ومكعب} طول ضلعه $3$؛ وكومة على هيئة L مصنوعة من لوح
$3 \times 2 \times 1$ وفوقه لوح $1 \times 2 \times 1$.
\end{exercise}

\begin{exercise}[$\star$]\label{exo:g5:solids:7}
يسع صندوق $12$ مكعبًا واحديًّا بالضبط في طبقة واحدة. أعطِ شكلين
ممكنين للطبقة ($L \times w$)، وأعطِ لكل منهما \omterm{def:g5:solids:volume}{حجم} الصندوق إن
كانت فيه $3$ طبقات كهذه.
\end{exercise}

\begin{exercise}[$\star$]\label{exo:g5:solids:8}
مع $27$ مكعبًا واحديًّا، أي \omterm{def:g5:solids:def}{مكعب} تستطيع أن تبني؟ ومع $64$؟ (فكّر في
عدّ الطبقات بالعكس.)
\end{exercise}

\begin{exercise}[$\star\star$]\label{exo:g5:solids:9}
مكعبات سكر طول ضلعها $1$ سم تأتي في صندوق طوله $10$ سم، وعرضه $5$ سم،
وارتفاعه $4$ سم (قياسات داخلية). فكم \omterm{def:g5:solids:def}{مكعب} سكر يدخل فيه؟ وإذا
استعملت أسرة $6$ مكعبات في اليوم، فكم يومًا يكفيها الصندوق
(قرِّب تقريبًا معقولًا)؟
\end{exercise}

\begin{exercise}[$\star\star$]\label{exo:g5:solids:10}
\omterm{def:g5:solids:def}{مكعب} $3 \times 3 \times 3$ دُهن بالأحمر من الخارج، ثم قُطع
إلى $27$ مكعبًا واحديًّا. فكم مكعبًا صغيرًا عليه دهان على
$3$ \omterm{def:g5:solids:def}{أوجه} بالضبط؟ وعلى $2$؟ وعلى $1$؟ وعلى لا شيء؟ (والتحقّق: \omterm{def:g1:addition:def}{مجموع}
الأعداد الأربعة يجب أن يكون $27$.)
\end{exercise}
